\subsection{Three photons of each helicity}\label{sec:eight-balanced}

For $\M(1_e,2_E,3^+,4^+,5^+,6^-,7^-,8^-)$ put
$(a,b,c)=(3,4,5)$ and $(h,i,j)=(6,7,8)$. Use the three-photon currents
$D$ of Eq.~\eqref{eq:eight-two-bridge}. Let $\Sigma$ sum the two triples
of labels independently over their six permutations. With the common
physical reference $p_1$, the result is
\begin{align}
 \M^{\phys}_{+++---}=-(\sqrt2 e)^6\,\Sigma\!
 \left\{\frac{D(abc|hij)+D(hij|abc)}{36}
      +\frac{D(abh|cij)+D(ahi|bcj)}{4}\right\}.
 \label{eq:eight-balanced-phys}
\end{align}
Each extreme partition occurs once. There are nine partitions with one
negative on the electron side and nine with two. Each partition contains
the $3!^2$ orders within its two currents.

For the spurious form choose leg 6 as the angle reference of each positive
photon. Use square references $7,8,6$ for negative legs $6,7,8$,
respectively. In the following sum first assign the six labels, then
look up each photon's reference in this fixed table. Reversing the
charged line exchanges the two complementary partitions with a minus
sign. The same reference table works on both sides, giving
\begin{equation}
 \M^{\spur}_{+++---}=-(\sqrt2 e)^6(1-P_{12})\,\Sigma\!
 \left\{\frac{D(abc|hij)}{36}+\frac{D(abh|cij)}4\right\}.
 \label{eq:eight-balanced-spur}
\end{equation}
The expression in braces generates ten partitions. The endpoint
exchange supplies their ten complements. Parity and a photon relabeling
map this helicity sector to itself. On 100,000 phase-space and spin samples
both expressions agreed with the signed sum of 720 diagrams to
$8.11\times10^{-11}$. On a separate 100,000 momenta the largest relative
differences from the polarized MadGraph rate were $4.45\times10^{-12}$ and
$8.57\times10^{-12}$ for the physical and spurious forms, respectively.

\subsection{Two fermion pairs and four photons}\label{sec:eight-two-pairs}

Pair $(1_e,2_E)$ and $(4_e,3_E)$ for the direct term. Define $K(L|R)$
by attaching the photons in $L$ and $R$ to the first and second pair.
Let $\mathcal P=(1-P_{14})(1-P_{23})\mathcal S_+\mathcal S_-$,
where each $\mathcal S$ sums all permutations of photons of that helicity
with unit weight. In each row $a,b,c,d$ denote successive positive
photons and $h,i$ successive negative ones, in external order. The
amplitudes are
\begin{align}
 \M_{++++} &= \mathcal P\left\{
  \frac{K(\varnothing|abcd)}{24}+\frac{K(a|bcd)}6+\frac{K(ab|cd)}8\right\},\nonumber\\
 \M_{+++-} &= \mathcal P\left\{
  \frac{K(\varnothing|abch)+K(h|abc)}6
   +\frac{K(a|bch)+K(ah|bc)}2\right\},\nonumber\\
 \M_{++--} &= \mathcal P\left\{
  \frac{K(\varnothing|abhi)+K(hi|ab)}4
  +\frac{K(h|abi)+K(a|bhi)+K(ah|bi)}2\right\}.
 \label{eq:eight-four-fermions}
\end{align}
For example, $K(ab|cd)$ is unchanged under either exchange within a
pair of photons and under exchange of the two complete charged lines.
These three factors of two give its divisor eight. In the mixed
$K(ah|bi)$ term only the joint line exchange survives.

The same spinor contraction serves both pole choices. For a pair $(t,u)$
write $J_{tu}(L)$ as a sum of angle-square spinor pairs. Contract
$(|x\rangle,|y])$ from one current with $(|v\rangle,|w])$ from the
other as $\langle xv\rangle[yw]$ and extend linearly. Then
\begin{equation}
 K(L|R)=-(\sqrt2 e)^6\,
 \frac{J_{12}(L)\cdot J_{43}(R)}
      {(p_1+p_2+\sum_{i\in L}p_i)^2}.
 \label{eq:eight-four-kernel}
\end{equation}
To construct $J_{tu}(L)$, sum over subsets $B\subseteq L$ and add the
two orientations
\[
 (W_t^\angle(B),W_u^\square(L\setminus B))
  +(W_u^\angle(L\setminus B),W_t^\square(B)).
\]
The two orientations give the ordinary massive vector current at
$L=\varnothing$. For a nonempty set $B$ write $Q_t(B)=p_t+\sum_{i\in B}p_i$
and $\eta_t=+1$ for the electron and $-1$ for the positron.
The wave on one endpoint satisfies
\begin{align}
 [Q_t(B)^2-m^2] W_t^\chi(B)
 &=\sum_{i\in B}\!\left[\eta_t\sqrt2\,
 Q_t(B\setminus i)\cdot\varepsilon_i+F_{ti}^\chi\right]W_t^\chi(B\setminus i)
 \nonumber\\*[-2pt]
 &\quad-\sum_{\{i,j\}\subseteq B}
 (\varepsilon_i\cdot\varepsilon_j)W_t^\chi(B\setminus\{i,j\}),
 \qquad W_t^\chi(\varnothing)=|\mathbf t\rangle_\chi .
 \label{eq:eight-charged-wave}
\end{align}
Here $B\setminus i$ means removal of that photon. On a square spinor
$|x]$, a positive photon gives $F_{ti}^\square|x]=\eta_t[i x]|i]$.
On an angle spinor $|x\rangle$, a negative photon gives
$F_{ti}^\angle|x\rangle=\eta_t\langle i x\rangle|i\rangle$. The other
actions vanish. The two-photon term follows by commuting the momentum
through the two adjacent polarization matrices and summing their
orders. Thus each subset wave can be reused when a fourth photon is
added.

There is also a shorter summand. For a line $(t,u)$ carrying the set
$L$, keep only the first orientation in $J_{tu}(L)$, with the angle
spinor on the electron. Apply $1+(-1)^{|L|}P_{tu}$ to this chiral
current. Exchanging the endpoints interchanges the two orientations.
Each single-photon insertion in Eq.~\eqref{eq:eight-charged-wave}
changes sign, and the contact that absorbs two photons is even. This
gives the indicated sign at every order. Use the two endpoint exchanges
independently before the outer permutation sum. Only one angle-square
contraction per pair of subset waves then has to be written.

For AllPhysical take the electron momentum of the emitting pair as
each polarization reference. Use its projected spinors as in
Sec.~\ref{sec:eight-two-minus}, with $p_1$ replaced by this momentum.
For AllSpurious use the first photon of opposite helicity if one
exists. If all helicities agree, use the next photon cyclically.
A line current contains every attachment order, so its Ward identity
allows these choices for each fixed partition. Apply the permutation
operators to the complete kernel including its references.
Parity gives the remaining helicities.

\subsection{Three fermion pairs and two photons}\label{sec:eight-three-pairs}

Use the direct pairs $A=(1,2)$, $B=(3,4)$ and $C=(5,6)$ with an
electron first in each. The exchanged photons connect the path $A,B,C$.
Write $H(L|M|R)$ for external photon sets on the left, middle and right
lines, respectively. Write $0$ for an empty set in the following two
equations and let $\mathcal A$ antisymmetrize the three electron and the
three positron labels independently without normalization. For $i=7,j=8$,
\begin{align}
 \M_{++} &=\mathcal A(1+P_{ij})\left\{
 \frac{H(ij|0|0)}2+\frac{H(0|ij|0)}4
 +\frac{H(i|0|j)}2+H(i|j|0)\right\},\nonumber\\
 \M_{+-} &=\mathcal A\left\{
 H(ij|0|0)+\frac{H(0|ij|0)}2+H(i|0|j)
 +H(i|j|0)+H(j|i|0)\right\}.
 \label{eq:eight-six-fermions}
\end{align}
For opposite helicities, emission of the positive photon from the middle
and emission of the negative photon from the middle require separate terms.
For two equal helicities they are exchanged by $P_{ij}$.

The following open-line tensor specifies each kernel, including emission
between two exchanged photons. At each attachment use $G_i=\not\!\varepsilon_i$
for an external photon or $G_i=\gamma^{\mu_i}$ for an exchanged one.
For a charged pair $(t,u)$ in electron, positron order, define
\begin{equation}
 T_{tu}(1,\ldots,k)=\sum_{\pi\in S_k}
 \bar v_u G_{\pi_1}
 \prod_{j=1}^{k-1}
 \left[\frac{m-\not\! Q_j}{Q_j^2-m^2}\,G_{\pi_{j+1}}\right]u_t,
 \qquad Q_j=p_u+\sum_{b=1}^{j}q_{\pi_b}.
 \label{eq:eight-ordered-tensor}
\end{equation}
The attachment numbers in $T$ label its arguments locally. For an external
photon $q_i=p_i$. For an exchanged photon $q_i$ is the sum of the external
momenta on its far side as seen from the charged line in question. The two
lines at that photon therefore use opposite momenta. Each ordered product
runs from the positron bra to the electron ket. With
$k_A=p_1+p_2+\sum_{i\in L}p_i$ and $k_C=p_5+p_6+\sum_{i\in R}p_i$,
\begin{equation}
 H(L|M|R)=-e^6\frac{
 T_{12}^{\mu}(L)\,T_{34,\mu\nu}(M)\,T_{56}^{\nu}(R)}
 {k_A^2 k_C^2}.
 \label{eq:eight-six-kernel}
\end{equation}
The displayed Lorentz indices denote the exchanged attachments in
Eq.~\eqref{eq:eight-ordered-tensor}. The physical reference is the
electron on the line emitting the photon. In the spurious form the two
external photons use each other as their references.

One can shorten $H$ in the same way. In either terminal line keep the
angle component immediately before its exchanged photon in the ordered
word of Eq.~\eqref{eq:eight-ordered-tensor}. In the middle line keep
only orders with the exchange attached to $A$ before the exchange
attached to $C$, reading from the positron bra. Call this kernel $H_0$.
Then, for each term in Eq.~\eqref{eq:eight-six-fermions},
\begin{align}
 H(L|M|R)={}&[1+(-1)^{|L|}P_{12}]
 [1-(-1)^{|M|}P_{34}]
 [1+(-1)^{|R|}P_{56}]\,H_0(L|M|R).
 \label{eq:eight-six-ends}
\end{align}
Each terminal exchange restores the other chirality. Reversing the
middle line supplies the opposite order of the two virtual photons.
The extra minus sign there follows from the second vector vertex.
All positions of the real photons relative to the virtual ones are
included in each retained order.

\subsection{Four fermion pairs}\label{sec:eight-four-pairs}

With no external photons, both pole choices give the same expression.
There are two ways to connect four charged lines with three photon
exchanges. In the star one line has three vector attachments and in
the path two lines have two. Pair the legs as $(1,2),(3,4),(5,6),(7,8)$,
and let $\mathcal A$ antisymmetrize $1,3,5,7$ and $2,4,6,8$
independently. Let $U^{\mu\nu\cdots}_{tu}$ be the single ordered product in
Eq.~\eqref{eq:eight-ordered-tensor}, with the attachments in the displayed
index order. Write $J_{tu}^{\mu}=\bar v_u\gamma^\mu u_t$. The amplitude is
\begin{align}
 \M=e^6\mathcal A\!\left[
 \frac{U_{12}^{\mu\nu\rho}J_{34,\mu}J_{56,\nu}J_{78,\rho}}
      {s_{34}s_{56}s_{78}}
 +\frac{(1-P_{34})(1-P_{56})}{2}
 \frac{J_{12}^{\mu}U_{34,\mu\nu}U_{56}^{\nu\rho}J_{78,\rho}}
      {s_{12}s_{1234}s_{78}}\right].
 \label{eq:eight-eight-fermions}
\end{align}
The three star leaves can occur in any order. Every permutation of
these leaves exchanges electrons and positrons together, so its sign
in $\mathcal A$ is positive. The six permutations supply the six
distinct orders on the central line. On the path, reversing either
two-vector line exchanges its two vector attachments with a minus
sign. The two factors $1-P$ therefore supply the four attachment
orders. The divisor two accounts for the two orientations of the path.
There are $4!$ matchings. For each matching the four choices of star
center give $4\cdot 3!$ diagrams, and the twelve paths give $12\cdot 2^2$
more. The sum accounts for all 1,728 tree diagrams.

The same two terms have a short expression in spinor brackets. Contract
a terminal pair $(c_e,d_E)$ into its neighbour, keeping one chirality
of the pair. The other chirality is supplied by $1+P_{cd}$. Suppress the
factor two from this Fierz contraction and denote the resulting endpoint
by a tilde. For a positron bra $b$ or an electron ket $t$, use
$Q_x=p_x+p_c+p_d$ and the following four spinors:
\begin{align}
 \langle\widetilde b|&=\langle d b\rangle[c|Q_b-m[b c]\langle d|,&
 [\widetilde b|&=-[b c]\langle d|Q_b+m\langle d b\rangle[c|,\nonumber\\
 |\widetilde t\rangle&=\langle d t\rangle Q_t|c]-m[c t]|d\rangle,&
 |\widetilde t]&=[c t]Q_t|d\rangle-m\langle d t\rangle|c].
 \label{eq:eight-dressed-leaf}
\end{align}
For the star use pair $(3,4)$ for $\widetilde2$ and $(7,8)$ for
$\widetilde1$. For the path use $(1,2)$ for $\widetilde4$ and $(7,8)$
for $\widetilde5$. The two scalar numerators are
\begin{align}
 N_{\rm star}&=\langle6\,\widetilde1\rangle[\widetilde2\,5]
                 +\langle\widetilde2\,6\rangle[5\,\widetilde1],\nonumber\\
 N_{\rm path}&=\langle63\rangle[\widetilde4\,\widetilde5]
           -\langle\widetilde4\,6\rangle[3\,\widetilde5]
           -\langle3\,\widetilde5\rangle[\widetilde4\,6]
           +\langle\widetilde4\,\widetilde5\rangle[63].
 \label{eq:eight-star-path-numerators}
\end{align}
With $S_{ij}=1+P_{ij}$ and $A_{ij}=1-P_{ij}$, the bracket formula reads
\begin{align}
 \M=8e^6\mathcal A\!\left[
 S_{34}S_{56}S_{78}\frac{N_{\rm star}}{s_{34}s_{56}s_{78}\tau_{234}\tau_{178}}
 +\frac{A_{34}A_{56}S_{12}S_{78}}2\,
 \frac{N_{\rm path}}{s_{12}s_{1234}s_{78}\tau_{124}\tau_{578}}\right].
 \label{eq:eight-eight-weyl}
\end{align}
Here the local exchanges act inside the outer fermion antisymmetrization.
The factor eight comes from the three photon contractions. This form
exhibits all five physical propagators in each term.

For numerical evaluation, the middle photon in the path divides the legs
into two sets, each containing two electrons and two positrons. There are
$\binom42^2/2=18$ unordered divisions. For each division, first sum the
signed exchanges within each four-leg current, then contract the two
currents across that photon. On the star, choosing the middle pair and
the two sets of three dressed legs leaves $16\cdot9=144$ assignments.
Antisymmetrizing the two equal-charge legs within each triple before
joining the waves supplies the other four images per assignment. This
evaluation order retains the relative sign between the star and the path.


At eight points the four families have $2^6$, $2^4$, $2^2$ and one
labeled photon-helicity assignment, respectively. Photon permutations
and parity reduce these 85 assignments to ten representatives, four
with one charged pair, three with two pairs, two with three pairs and
one with four pairs. The massive external spins remain open in every
formula.
